% Generated by tools/edn_latex.py from reports/edn.md.
% Do not edit: change reports/edn.md and run `make edn`.
\ednentry{
  id = {EDN-20},
  title = {The DagsHub remote stays public},
  date = {2026-09-29},
  milestone = {M2: Reproducibility},
  topic = {Data Versioning, Privacy},
  participants = {@lukas2510 (Scrum Master). Pending confirmation by the full team at the next sprint review.},
  decision = {The DagsHub repository that serves as our DVC remote and MLflow tracking server stays public.
We accept that the raw file's PII columns are readable by anyone with a DagsHub account, which follows from EDN-25 tracking the raw file with \texttt{dvc add} and pushing it to that remote.
The additional exposure this creates is small, because the identical file is already publicly downloadable from the pinned Zenodo DOI under the MIT license.
It is not nothing: hosting our own copy makes the team a publisher of that personal data in its own right and under its own name, which the earlier publication mitigates but does not undo.},
  rationale = {\emph{Alternatives considered:}
\begin{itemize}[leftmargin=*, topsep=2pt]
\item \textbf{Option A (chosen): keep the repository public and say so.}
\newline Pros: graders and supervisors can inspect code, data and experiments without being added as collaborators, which is the reason the course puts the project on DagsHub in the first place; it matches how the source dataset is already published.
\newline Cons: the raw file, which still carries \texttt{vin}, \texttt{street}, \texttt{seller\_\allowbreak{}company\_\allowbreak{}name} and coordinates, can be pulled by every logged-in DagsHub user; re-hosting it is a publishing act of our own, so ``the author already published it'' reduces the marginal risk but does not transfer the responsibility; the decision has to be revisited the moment we host data that is not already public elsewhere.
\item \textbf{Option B: make the repository private.}
\newline Pros: the rationale first given for pushing the raw file would hold as written; the PII sits behind an access wall.
\newline Cons: every grader, supervisor and teammate then needs a manual collaborator entry; it buys little real protection, because the identical file stays publicly downloadable from Zenodo either way.
\end{itemize}
\par
\emph{Why:} The question only came up because the premise of the raw-data decision turned out to be false.
The first version of EDN-07, on the issue \#3 branch, justified pushing the raw file with ``the DagsHub repo is private (anonymous access is redirected to the login page and an anonymous data download returns 401)''.
Checked on 2026-09-29: DagsHub refuses anonymous access to everything, public repositories included.
A known-public control repository (\texttt{DAGsHub-\allowbreak{}Official/\allowbreak{}dagshub-\allowbreak{}docs}, \texttt{private: false}) answers anonymous requests exactly like ours, and so does a repository that does not exist, so the observation has no discriminating power.
The authenticated API reports \texttt{private: false} for both mirrors that existed at the time.
Given a real choice between hiding a file that is already public and saying openly that it is public, Lukas chose the second: the protection gained would be nominal, while the access cost for graders and supervisors would be real.
Under EDN-19 the organisation is the team's own, so the visibility is ours to change and this decision can be revisited at any time without asking anyone outside the team.
Recorded deliberately as an accepted risk rather than a solved problem, because the two are not the same thing: the marginal exposure is small, but we are still the ones publishing personal data, and an entry that claimed the concern was spent would not survive a reviewer who cares about privacy.
State after the decision: \texttt{mark.\allowbreak{}welf.\allowbreak{}atzberger/\allowbreak{}MLOPS\_\allowbreak{}Recommenditos}, the repository that held the raw copy, was deleted on 2026-09-29, and the object returns 404. The team therefore publishes no personal data on any DagsHub remote at present, and this decision governs whatever remote EDN-19 settles on rather than an exposure that exists today.},
  ai = {Information seeking, Alternative generation, Alternative assessment},
  response = {Used as input for further analysis},
  assessment = {AI found that the privacy premise of the raw-data decision was false, using a control repository to show that anonymous refusal does not distinguish public from private, and put both ways out to Lukas without recommending either, since the trade-off is about how open the team wants to be rather than a technical question. Lukas decided to keep the repository public. The finding is what changed the outcome here; the decision itself was not AI's to make.},
  aievidence = {Claude Code session on 2026-09-29, prompts translated from German: after the finding was presented with the two options (``either mark switches it to private ... or you keep it public and write that honestly into the raw-data entry''), Lukas answered ``we keep it public''.},
  evidence = {EDN-19, EDN-25, \href{https://github.com/mlops-2627q1-mds-upc/MLOPS_Recommenditos/pull/26}{PR \#26}, \href{https://github.com/mlops-2627q1-mds-upc/MLOPS_Recommenditos/pull/27}{PR \#27}.},
}
